\section{Arquitectura de Persistencia, CDC y Plataforma Web}
\label{sec:arquitectura_datos}

Para superar las limitaciones inherentes a los archivos planos y permitir la interoperabilidad transaccional multiusuario, el proyecto evolucionó desde una base de datos local embebida hacia una arquitectura distribuida basada en \emph{Neon Serverless PostgreSQL} (compatible con versiones 17.5 y 18.x), alojada en la nube con conexión cifrada SSL obligatoria (\texttt{sslmode=require}).

\subsection{Esquema Físico Relacional y Modelado de Datos}

El modelo relacional está compuesto por 16 tablas organizadas en cuatro subsistemas lógicos:
\begin{enumerate}[label=\alph*)]
  \item \textbf{Catálogo Académico Institucional:} Tablas \texttt{docentes}, \texttt{materias} y \texttt{oferta\_}\allowbreak\texttt{academica\_}\allowbreak\texttt{semestral}. Modela la estructura curricular de la Carrera de Informática de la UMSA y gestiona la evolución temporal de asignaturas y paralelos mediante una estrategia de dimensión de cambio lento (SCD Tipo 2).
  \item \textbf{Tabla de Hechos Central:} La tabla \texttt{mensajes} almacena los 23.371 mensajes del corpus, desnormalizando atributos calculados de alto acceso (polaridad léxica, emoción dominante, banderas multimodales para OCR, stickers y GIFs, e índices numéricos de intensidad afectiva).
  \item \textbf{Matrices Analíticas de Alta Densidad:} Las tablas \texttt{tfidf\_mensajes} (122.971 registros ponderados), \texttt{ngrams\_frecuencias} (27.334 filas) y \texttt{metricas\_temporales} desacoplan el cálculo computacional intensivo de la capa de visualización.
  \item \textbf{Trazabilidad y Menciones Académicas:} La tabla \texttt{menciones\_docentes} vincula los mensajes con las entidades docentes identificadas mediante expresiones regulares compiladas (\texttt{alias\_regex}).
\end{enumerate}

\begin{table}[htbp]
  \centering
  \small
  \caption{Métricas Operativas Globales y Composición Multimodal del Corpus de la UMSA}
  \label{tab:kpis_corpus}
  \resizebox{\textwidth}{!}{
  \begin{tabular}{lrr}
    \toprule
    Métrica / Dimensión & Valor Observado & Observación Metodológica \\
    \midrule
    Total de mensajes persistidos          & 23.371 & Corpus histórico consolidado en Neon PostgreSQL \\
    Usuarios únicos remitentes             & 2.051  & Estudiantes, auxiliares y coordinadores \\
    Promedio de mensajes por usuario       & 11,4   & Distribución de cola larga con alta dispersión \\
    Total de elementos multimodales        & 4.919  & Imágenes, stickers y animaciones GIF \\
    Imágenes procesadas con OCR            & 1.782  & Capturas de comunicados y boletas de inscripción \\
    Stickers analizados                    & 2.180  & Clasificados por carga afectiva y emoji asociado \\
    Animaciones GIF documentadas           & 24     & Reacciones animadas catalogadas \\
    Reacciones interactivas de Telegram    & 2.575  & Emojis directos en globos de mensaje \\
    Modismos y jergas locales de la UMSA   & 538    & Catálogo léxico contextual paceño \\
    Términos procesados en TF-IDF          & 122.971 & Matriz dispersa de términos ponderados \\
    \bottomrule
  \end{tabular}
  }
  \vspace{4pt}
  \begin{minipage}{0.94\textwidth}
    \footnotesize
    \textit{Nota.} Datos extraídos en tiempo real desde las tablas operativas \texttt{mensajes}, \texttt{menciones\_docentes} y \texttt{tfidf\_mensajes} en Neon Serverless PostgreSQL (región AWS us-east-2 / sa-east-1).
  \end{minipage}
\end{table}


\subsection{Captura de Cambios de Datos en Tiempo Real (CDC)}

La trazabilidad del sistema se garantiza a nivel de motor de base de datos mediante la tabla \texttt{bitacora\_cambios} y una función disparadora genérica (\texttt{fn\_registrar\_bitacora\_cambio()}) escrita en PL/pgSQL (Listado~\ref{lst:cdc_trigger}). Ante cualquier operación de actualización (\texttt{UPDATE}) o eliminación (\texttt{DELETE}), el trigger intercepta los registros \texttt{OLD} y \texttt{NEW}, compara campo por campo utilizando funciones de introspección JSONB y genera un delta granular (\texttt{diff\_modificaciones}).

\begin{lstlisting}[language=SQL, caption={Funcion trigger PL/pgSQL para captura diferencial de cambios (CDC) en formato JSONB}, label={lst:cdc_trigger}]
CREATE OR REPLACE FUNCTION fn_registrar_bitacora_cambio()
RETURNS TRIGGER AS $$
DECLARE
    v_old_json JSONB := NULL;
    v_new_json JSONB := NULL;
    v_diff JSONB := '{}'::JSONB;
    v_key TEXT;
    v_pk_val TEXT := NULL;
BEGIN
    IF (TG_OP = 'UPDATE') THEN
        v_old_json := to_jsonb(OLD);
        v_new_json := to_jsonb(NEW);
        v_pk_val := COALESCE(v_new_json->>'message_id', v_new_json->>'id');

        -- Deteccion granular de atributos modificados
        FOR v_key IN SELECT jsonb_object_keys(v_new_json)
        LOOP
            IF v_new_json->v_key IS DISTINCT FROM v_old_json->v_key THEN
                v_diff := v_diff || jsonb_build_object(
                    v_key, jsonb_build_object('anterior', v_old_json->v_key, 'nuevo', v_new_json->v_key)
                );
            END IF;
        END LOOP;

        IF v_diff <> '{}'::JSONB THEN
            INSERT INTO bitacora_cambios (
                tabla_afectada, operacion, registro_id, datos_previos,
                datos_nuevos, diff_modificaciones, fecha_hora
            ) VALUES (
                TG_TABLE_NAME, TG_OP, v_pk_val, v_old_json, v_new_json, v_diff, clock_timestamp()
            );
        END IF;
        RETURN NEW;
    END IF;
    RETURN NULL;
END;
$$ LANGUAGE plpgsql SECURITY DEFINER;
\end{lstlisting}


\subsubsection{Optimización y Prevención de Amplificación de Escritura}
Durante la fase experimental, la inserción masiva de 12.581 tokens y registros lingüísticos generó un fenómeno de amplificación de escritura (\emph{write amplification}), saturando la bitácora con más de 12.000 eventos atómicos. Como solución ingenieril documentada en \path{database/optimize_bitacora.py}, se desactivaron los triggers de fila sobre tablas analíticas transitorias (\texttt{tokens\_frecuencia} y \texttt{shadow\_linguistics}), canalizando la trazabilidad de lotes hacia \texttt{bitacora\_procesos\_mineria}.

\subsection{Plataforma Web y Analítica Visual en Tiempo Real}

El consumo de datos y la interacción del usuario se centralizan en una aplicación web modular desarrollada en Python mediante el microframework Flask, organizada bajo el patrón \emph{Application Factory} con dos Blueprints:
\begin{itemize}
  \item \textbf{Blueprint Web (\texttt{app/routes/web}):} Renderizado del lado del servidor (SSR) mediante plantillas Jinja2, diseño responsive basado en tokens de Material Design 3, modo oscuro nativo anti-FOUC y visualizadores especializados.
  \item \textbf{Blueprint API REST v1 (\texttt{app/routes/api\_v1}):} Endpoints de telemetría y el controlador de ingesta (\texttt{/api/v1/ingest/upload}).
\end{itemize}

\subsubsection{Aislamiento In-Memory para Ingesta Multimodal}
Para proteger la integridad del corpus oficial de la UMSA, la subida de nuevos paquetes de chat (archivos JSON o archivos comprimidos ZIP de Telegram Desktop con imágenes y stickers) opera bajo un esquema de \emph{aislamiento de escritura cero} (\emph{zero-write isolation}). Tal como se aprecia en el Listado~\ref{lst:ingest_memory}, el motor en Python (\texttt{app.core.emotica}) parsea, normaliza y puntúa las emociones íntegramente en memoria RAM, realizando únicamente consultas de solo lectura (\texttt{SELECT}) a Neon PostgreSQL para contrastar el nuevo chat contra el promedio histórico.

\begin{lstlisting}[language=Python, caption={Ingesta en memoria con aislamiento de transacciones de escritura (\texttt{ingest.py})}, label={lst:ingest_memory}]
@api_v1_bp.route("/ingest/upload", methods=["POST"])
def ingest_upload_view():
    """
    Procesa lotes de Telegram (JSON o ZIP) 100% en memoria en Python.
    Compara las metricas resultantes contra el baseline de Neon (Read-Only).
    """
    archivos = request.files.getlist("files")
    if not archivos:
        return jsonify({"error": "No se recibieron archivos"}), 400

    # 1. Parseo y extraccion en memoria sin mutacion de base de datos
    df_mensajes = parsear_paquete_telegram(archivos)
    
    # 2. Preprocesamiento y scoring lexico
    df_limpio = preprocess_chat(df_mensajes)
    scores_df = score_emotions(df_limpio)
    
    # 3. Consulta de solo lectura (SELECT) contra Neon PostgreSQL
    baseline = consultar_baseline_corpus_neon()
    timeline = calcular_timeline_chat(df_limpio)

    return jsonify({
        "status": "success",
        "total_mensajes": len(df_mensajes),
        "timeline": timeline,
        "comparativa_neon": {"chat": scores_df, "baseline": baseline}
    }), 200
\end{lstlisting}


\subsubsection{Ingeniería de Motores Gráficos Especializados}
A diferencia de tableros convencionales que delegan la representación a bibliotecas genéricas rígidas, la plataforma incorpora dos motores visuales avanzados:
\begin{enumerate}
  \item \textbf{Motor Spline SVG Nativo (\texttt{timeline\_chart.js}):} Renderiza curvas continuas de Bézier sobre el DOM mediante elementos vectoriales \texttt{<path>}, soportando selector dinámico de granularidad (anual, semestral, mensual, 24 horas), crosshair magnético con interpolación de coordenadas y tooltips glassmórficos.
  \item \textbf{Flujo ETL en Sankey (\texttt{sankey\_engine.js}):} Emplea Apache ECharts para graficar la reducción de dimensionalidad y pérdida de volumen del dato desde su estado crudo hasta los tokens afectivos útiles.
\end{enumerate}
